\section{Literature Review}

Large language models (LLMs) and transformer-based architectures have significantly changed natural language processing by enabling contextual representation learning, large-scale text generation, and flexible adaptation to downstream tasks. The Transformer architecture introduced self-attention as the core mechanism for sequence modeling, removing the need for recurrent or convolutional structures and enabling more scalable parallel training \cite{vaswani2017attention}. Building on this architecture, BERT introduced bidirectional contextual pretraining and demonstrated strong performance across several natural language understanding tasks \cite{devlin2019bert}. These developments established the foundation for modern LLMs and domain-specific language models. However, when LLMs are used in knowledge-intensive settings, their responses may be fluent but unsupported, outdated, or difficult to verify. This limitation is especially important in biomedical and healthcare applications, where factual accuracy and source traceability are essential.

Retrieval-augmented generation (RAG) was introduced to address some of these limitations by combining parametric knowledge stored in a pretrained generation model with non-parametric knowledge retrieved from external documents \cite{lewis2020rag}. In a RAG system, retrieved passages are provided as context to the generator, allowing the model to produce responses that are more closely connected to source evidence. Earlier retrieval-augmented language modeling work, such as REALM, showed that neural models can retrieve external documents during pretraining and downstream question answering, thereby improving open-domain question answering and providing a more interpretable link between generated answers and retrieved evidence \cite{guu2020realm}. Similarly, retrieval-enhanced language models such as RETRO demonstrated that large-scale retrieval from external corpora can improve language modeling efficiency and knowledge-intensive task performance \cite{borgeaud2022retro}. These studies show that retrieval can serve as an explicit memory mechanism for language models.

Open-domain question answering research has also shown that retrieval quality directly affects answer quality. Dense Passage Retrieval (DPR) demonstrated that dense dual-encoder representations can retrieve relevant passages effectively for open-domain question answering and outperform traditional sparse retrieval baselines in several settings \cite{karpukhin2020dpr}. Unsupervised dense retrieval methods have also shown that contrastive learning can improve retrieval representations without relying entirely on supervised relevance labels~\cite{izacard2022contriever}. Fusion-in-Decoder (FiD) further showed that generative models can benefit from conditioning on multiple retrieved passages, allowing the decoder to integrate evidence across several retrieved contexts \cite{izacard2021fid}. These approaches highlight the importance of both retrieving relevant evidence and designing generation models that can effectively use retrieved passages.

Traditional lexical retrieval remains important in RAG pipelines because exact term matching is often useful when queries contain technical terminology, named entities, abbreviations, or domain-specific phrases. BM25 and related probabilistic relevance models remain widely used because they provide efficient and interpretable keyword-based retrieval \cite{robertson2009bm25}. However, lexical retrieval may miss semantically relevant passages when the query and document use different wording. Dense retrieval addresses this vocabulary mismatch by representing queries and documents as vectors in a shared semantic space \cite{karpukhin2020dpr}. Late-interaction retrieval models such as ColBERT and ColBERTv2 attempt to balance efficiency and semantic expressiveness by preserving token-level interactions while enabling scalable retrieval \cite{khattab2020colbert,santhanam2022colbertv2}. Sparse neural retrieval models such as SPLADE also attempt to combine lexical interpretability with learned expansion, producing sparse representations that can improve matching beyond exact surface terms \cite{formal2021splade}.

Because sparse and dense retrieval methods have complementary strengths, hybrid retrieval has become increasingly important for RAG systems. Benchmarks such as BEIR have shown that retrieval models must be evaluated across heterogeneous datasets because performance in one retrieval setting may not generalize to another \cite{thakur2021beir}. MS MARCO has also played an important role in training and evaluating large-scale passage retrieval systems for question answering and ranking tasks \cite{bajaj2016msmarco}. Hybrid retrieval combines lexical and semantic signals to improve retrieval coverage, especially when relevant evidence may be found through exact terminology or broader semantic similarity. OpenSearch supports hybrid search by combining keyword-based and vector-based retrieval signals, including score normalization and combination during query processing \cite{opensearchhybrid}. For the present study, this is directly relevant because the proposed system uses hybrid retrieval before reranking and answer generation.

Reranking is commonly used as a second-stage retrieval refinement method. In many RAG pipelines, the initial retriever returns a candidate set of passages, and a reranker then reorders those passages based on their relevance to the user query. This two-stage design is useful because the first-stage retriever is optimized for broad candidate recall, while the reranker can perform more focused relevance estimation. Cohere rerank models are designed to sort candidate texts according to semantic relevance to a specified query \cite{coherererank}. Reranking can therefore improve the quality of the final evidence context provided to the generator, particularly when the initial retrieval output contains partially relevant or weakly ranked chunks.

Biomedical and clinical NLP introduce additional challenges because medical language contains specialized terminology, abbreviations, domain-specific relationships, and context-dependent meanings. BioBERT showed that continued pretraining on biomedical corpora such as PubMed abstracts and PMC full-text articles improves performance on biomedical named entity recognition, relation extraction, and question answering \cite{lee2020biobert}. PubMedBERT further demonstrated that pretraining a language model from scratch on biomedical text can provide strong performance across biomedical NLP benchmarks \cite{gu2021pubmedbert}. ClinicalBERT adapted transformer-based representation learning to clinical notes and showed utility for modeling electronic health record text and predicting hospital readmission \cite{huang2019clinicalbert}. More recent large-scale clinical language models, such as GatorTron and GatorTronGPT, further demonstrate that scaling language models on clinical and biomedical text can improve clinical NLP and healthcare text generation tasks \cite{yang2022gatortron,peng2023gatortrongpt}.

Biomedical question answering and clinical evidence retrieval require careful evaluation because generated responses may affect scientific interpretation or healthcare decision-making. Datasets and challenges such as BioASQ and PubMedQA have helped advance biomedical semantic indexing and biomedical question answering by requiring systems to retrieve and reason over biomedical literature \cite{tsatsaronis2015bioasq,jin2019pubmedqa}. MIMIC-IV provides a large-scale, deidentified electronic health record database that has supported a wide range of clinical informatics and machine learning research \cite{johnson2023mimiciv}. The PhysioNet release of MIMIC-IV version 3.1 further provides a citable database version for reproducible research workflows \cite{johnson2024mimiciv}. Although the present study focuses on document-grounded RAG rather than direct clinical prediction, these resources demonstrate the importance of reproducible datasets and evidence-based evaluation in biomedical informatics. Recent related studies have examined selective fine-tuning and GPT-based architectures for clinical text classification using EHR notes \cite{irany2026selective,irany2026generative}, while complementary educational research has investigated the impact of AI on student performance in engineering technology courses \cite{hossain2025impact}.

Recent studies on LLMs in medicine show both the promise and risk of generative AI in healthcare. Thirunavukarasu et al. reviewed applications of LLMs in medicine and emphasized the need for validation, safety assessment, and careful deployment in clinical settings \cite{thirunavukarasu2023llmmedicine}. Singhal et al. evaluated large language models on medical question answering and proposed human evaluation dimensions such as factuality, comprehension, reasoning, possible harm, and bias \cite{singhal2023medicalknowledge}. Ayers et al. compared physician and chatbot responses to patient questions and showed that generative systems may produce responses perceived as high quality and empathetic, while also highlighting the need for careful evaluation before clinical use \cite{ayers2023chatbot}. These studies reinforce the need for grounded, verifiable, and transparent LLM systems in healthcare-related information retrieval.

Hallucination remains a major limitation of neural text generation. Ji et al. define hallucination as generated content that is unsupported by or inconsistent with the source input or factual knowledge, and they identify hallucination as a persistent problem across natural language generation tasks \cite{ji2023hallucination}. In RAG systems, retrieval can reduce hallucination by supplying external evidence, but retrieval alone does not guarantee that every generated statement is supported. A model may still misinterpret retrieved passages, overgeneralize from incomplete evidence, or generate plausible claims that are not directly grounded in the provided context. Therefore, RAG systems require not only strong retrieval and reranking but also systematic grounding evaluation.

Claim verification research provides a useful foundation for evaluating grounded generation. FEVER introduced a large-scale fact verification dataset in which claims are classified as supported, refuted, or not enough information based on textual evidence \cite{thorne2018fever}. SciFact extended claim verification to scientific literature by requiring systems to identify evidence-containing abstracts and rationales for scientific claims \cite{wadden2020scifact}. These studies are relevant to citation-aware RAG because generated responses can be decomposed into factual claims and each claim can be evaluated against retrieved evidence. In biomedical and scientific settings, this approach is more informative than evaluating the generated answer only as a single response.

Recent evaluation frameworks further support fine-grained assessment of RAG outputs. FActScore decomposes long-form generated text into atomic facts and measures the percentage of facts supported by a reliable knowledge source \cite{min2023factscore}. ALCE evaluates LLM-generated answers with citations and considers dimensions such as fluency, correctness, and citation quality \cite{gao2023alce}. RAGAS provides a reference-free evaluation framework for RAG pipelines, including evaluation dimensions related to retrieval and generation quality \cite{es2024ragas}. These frameworks support the motivation for claim-level grounding evaluation in the present study. Instead of treating a generated answer as entirely correct or incorrect, claim-level evaluation allows each factual sentence to be assessed against the retrieved evidence.

Cloud-based RAG infrastructure has made it easier to build reproducible retrieval, embedding, indexing, and generation pipelines. Amazon Bedrock Knowledge Bases support RAG workflows by connecting foundation models to user-provided data sources and retrieving relevant information to support response generation \cite{awsbedrockkb}. Amazon Bedrock Knowledge Bases can manage document ingestion, chunking, embedding, and vector store integration, which reduces the engineering burden of constructing a RAG pipeline from individual components \cite{awsbedrockworkflow}. Amazon OpenSearch Serverless provides vector search capabilities for scalable similarity search and generative AI applications \cite{opensearchserverless}. In addition, OpenSearch hybrid search enables the combination of keyword and semantic search signals \cite{opensearchhybrid}. These services are directly relevant to this study because the proposed framework uses Amazon S3 for document storage, Amazon Bedrock Knowledge Bases for ingestion and retrieval, Titan Text Embeddings V2 for embedding generation, OpenSearch Serverless for vector storage, hybrid retrieval for candidate selection, and Cohere reranking for evidence prioritization.

Overall, the literature shows that RAG can improve factual grounding by connecting LLMs to external evidence, but reliable RAG requires careful design across retrieval, reranking, generation, and evaluation. Prior research has established the importance of dense retrieval, lexical retrieval, hybrid retrieval, reranking, biomedical language modeling, hallucination reduction, and claim-level evaluation. However, there remains a need for practical and reproducible frameworks that integrate managed RAG infrastructure with hybrid retrieval, reranking, conservative answer generation, and systematic grounding assessment. The present study addresses this gap by proposing a hybrid retrieval and reranking framework for citation-aware RAG and evaluating generated responses using claim-level grounding accuracy.